% Source: https://artofproblemsolving.com/community/c3h1672611_the_water_bottle_problem
% Stable identity: aops-topic-1672611-q01
\textbf{Where the surface is.} The water surface is a horizontal plane. The cylinder is symmetric about its centre $O$ (the midpoint of the axis), and any plane through $O$ that does not meet the two bases cuts it into two congruent halves (the half-turn about $O$ swaps them). The horizontal plane through $O$ meets the axis at angle $90^\circ-30^\circ=60^\circ$; it reaches the wall at axial distances $4\tan30^\circ=\dfrac{4}{\sqrt3}\approx2.3$ cm from $O$ on either side, well inside the $10$ cm to each base. So it cuts off exactly half the volume, and the water surface is the section of the cylinder by this plane.

\begin{asy}
usepackage("amsmath");
size(6cm);
defaultpen(fontsize(10pt));
real t=30;pair ax=dir(90-t),nn=dir(-t);
pair BL=-4*nn,BR=4*nn,TR=4*nn+20*ax,TL=-4*nn+20*ax,Cc=10*ax;
pair Lw=extension(BL,TL,Cc,Cc+(1,0)),Rw=extension(BR,TR,Cc,Cc+(1,0));
fill(BL--BR--Rw--Lw--cycle,paleblue);
draw(BL--BR--TR--TL--cycle,linewidth(0.9));
draw(Lw--Rw,blue+linewidth(1));
draw((0,0)--10*ax,dashed);draw((0,0)--(0,6.8),dashed);
draw(arc((0,0),6,90-t,90));label("$30^\circ$",7.4*dir(90-t/2));
dot(Cc);label("$O$",Cc,N);
label("$8$",(BL+BR)/2+1.3*dir(-120));
label("$20$",(BR+TR)/2,SE);
label("$L$",Lw,W);label("$R$",Rw,E);
label("\mbox{}",(min(currentpicture,user=true)+max(currentpicture,user=true))/2,black); // keeps the text after the figure black when LaTeX typesets the labels inline
\end{asy}

\textbf{The ellipse.} In the side view (the plane containing the axis and the vertical), the surface is the segment $LR$. Its length is the width of the cylinder, $8$, divided by $\cos30^\circ$, since $LR$ makes an angle of $30^\circ$ with a cross-section perpendicular to the axis:
\[
LR=\dfrac{8}{\cos30^\circ}=\dfrac{16}{\sqrt3}.
\]
Perpendicular to the page, the surface is $8$ wide (a diameter of the cylinder). The section of a cylinder by a slanted plane is an ellipse with semi-axes $\dfrac{8}{\sqrt3}$ and $4$ (it is the circle of radius $4$ stretched by the factor $\dfrac{1}{\cos30^\circ}$ in one direction), so its area is
\[
\pi\cdot\dfrac{8}{\sqrt3}\cdot4=\dfrac{32\pi}{\sqrt3}=\boxed{\dfrac{32\sqrt3}{3}\pi}\ \text{cm}^2\approx58.0\text{ cm}^2 .
\]
